% Copy-ready manuscript fragment. It should replace duplicated assembly prose
% rather than be inserted in addition to it.

\begin{theorem}[Canonical atom normal form]
\label{theorem-canonical-atom-normal-form}
Let $F$ be a finite triple system without isolated points. Suppose that $F$
is linear, every hyperedge-node of $I(F)$ is incident with a bridge, and every
Berge cycle of $F$ has even length. Then $E(F)$ has a canonical partition into
subsystems of the following two types:
\begin{enumerate}
\item a single triple;
\item $J^+$, where $J$ is a finite $2$-connected bipartite graph.
\end{enumerate}
Two distinct subsystems meet in at most one point, and the incidence graph
between the subsystems and the points lying in at least two of them is a
forest. Conversely, every forest assembly of such subsystems by disjoint
unions and one-point amalgamations satisfies the three displayed intrinsic
conditions.
\end{theorem}

\begin{proof}
Let $L=I(F)$ and let $B$ be the set of bridges of $L$. A hyperedge-node has
degree three in $L$ and is incident with a member of $B$. Its degree in $L-B$
is therefore at most two. Residual degree one is impossible, because every
nonbridge edge lies on a cycle and a cycle enters and leaves a
hyperedge-node through two distinct incidences. Thus every hyperedge-node has
residual degree zero or two.

Let $C$ be a cyclic block of $L$. Every hyperedge-node in $C$ has degree two
in $C$. Suppressing these nodes gives a graph $J_C$ on the point-nodes of
$C$. The graph has neither loops nor parallel edges: the former are excluded
because a triple has three distinct points, and the latter by linearity of
$F$. Since $C$ is a subdivision of $J_C$, the graph $J_C$ is $2$-connected.
Moreover, cycles of $J_C$ correspond to Berge cycles contained in $C$.
Consequently, every cycle of $J_C$ is even, and $J_C$ is bipartite.

For a hyperedge-node $e\in C$, its third incidence is a bridge. Its point
$p_e$ does not lie in $C$, since otherwise that incidence, together with a
path in $C$, would lie on a cycle. The points $p_e$ are pairwise distinct:
if $p_e=p_f$, a path in $C$ from $e$ to $f$, completed through their common
point, would put both bridge incidences on a cycle. The hyperedges represented
in $C$ therefore form a subsystem isomorphic to $J_C^+$.

A hyperedge-node outside every cyclic block cannot have residual degree two,
for its two surviving incidences would lie on a cycle. It has residual degree
zero and contributes a single-triple subsystem. These subsystems partition
$E(F)$.

Two distinct subsystems cannot share two points. Otherwise, Levi paths within
the two subsystems between the shared points would contain a cycle crossing
the two subsystems, whereas every Levi cycle belongs to one cyclic block. The
same argument shows that the subsystem--shared-point incidence graph is
acyclic: a cycle in that incidence graph expands to a closed Levi walk crossing
subsystem boundaries, and deleting spurs produces a Levi cycle with the same
property. It is therefore a forest. The cyclic blocks and the remaining
all-bridge hyperedges are canonical, and suppressing the degree-two
hyperedge-nodes determines each $J_C$ up to isomorphism.

Conversely, every listed subsystem is linear, every hyperedge-node has its
private incidence as a bridge, and every Berge cycle comes from a cycle of a
bipartite core. Forest assembly preserves linearity, localizes Berge cycles in
one subsystem, and preserves the private bridges. Hence the assembled system
satisfies the intrinsic conditions.
\end{proof}

\begin{corollary}[Minimal generators]
\label{corollary-minimal-obligatory-generators}
The class $\mathcal B$ is the smallest isomorphism-closed class containing
finite edgeless systems, one triple, and $J^+$ for every finite $2$-connected
bipartite graph $J$, and closed under finite disjoint unions and one-point
amalgamations.
\end{corollary}

\begin{proof}
Every listed generator belongs to the original definition of $\mathcal B$.
Conversely, apply Theorem~\ref{theorem-canonical-atom-normal-form} to the
isolated reduction and restore the isolated points by adjoining an edgeless
system.
\end{proof}

\begin{corollary}[One-point-indecomposable systems]
\label{corollary-indecomposable-obligatory-systems}
A finite reduced connected obligatory triple system with at least one
hyperedge is one-point indecomposable if and only if it is one triple or is
isomorphic to $J^+$ for a finite $2$-connected bipartite graph $J$.
\end{corollary}

\begin{proof}
If the canonical atom forest has more than one atom, a leaf atom separates at
its unique shared point. Thus an indecomposable system consists of one atom.
A single triple is indecomposable. Let $J$ be $2$-connected. Deleting a
private point from $I(J^+)$ removes only a leaf. Deleting a core point $v$
does not separate the hyperedge-nodes, because $J-v$ is connected and every
edge-node formerly incident with $v$ remains attached through its other core
endpoint. Hence $J^+$ admits no nontrivial one-point decomposition.
\end{proof}

\begin{remark}[Rooted-theta form of the bridge condition]
Let a hyperedge-node $e$ have point-neighbours $x,y,z$. None of the three
incidences at $e$ is a bridge if and only if $x,y,z$ lie in one component of
$I(F)-e$; equivalently, $I(F)$ contains a theta subgraph with branch vertex
$e$ whose three branches begin with the three incidences at $e$. Thus the
intrinsic bridge condition excludes a rooted theta using all incidences of a
hyperedge-node.
\end{remark}
